\section[Rédaction Bourbaki n\textsuperscript{o}\,288 : la loi exponentielle]{Rédaction Bourbaki n\textsuperscript{o}\,288 : la loi exponentielle pour la convergence compacte\protect\verifie{Bourbaki288}}\label{sec:bourbaki-288}

Cette section ne porte pas sur une cote du fonds de Montpellier. La rédaction n\textsuperscript{o}\,288 des archives Bourbaki, datée de 1958 (Archives Henri Poincaré, \href{https://archives-bourbaki.ahp-numerique.fr/items/show/689}{item 689}), est un contre-projet pour le chapitre~X de \emph{Topologie générale}, qui traite des espaces d'applications continues. Elle est placée dans cette partie parce que la topologie de la convergence compacte sur ces espaces est l'un des outils de base de l'analyse fonctionnelle des années 1950. Sa transcription reste locale pour l'instant ; on n'en cite ici que quelques mots.

\emph{La loi exponentielle.} Une fonction de deux variables $f(x, y)$ peut se lire comme une famille de fonctions d'une variable : à chaque $x$, on associe la fonction partielle $y \mapsto f(x, y)$. Pour les ensembles, ce passage est une bijection entre les applications $E \times F \to G$ et les applications de $E$ dans l'ensemble des applications $F \to G$. Pour les cardinaux, il s'écrit $|G|^{|E|\cdot|F|} = \bigl(|G|^{|F|}\bigr)^{|E|}$, d'où le nom de loi exponentielle ; les informaticiens parlent de curryfication. En topologie, on demande davantage. On note $C(E, G)$ l'espace des applications continues de $E$ dans $G$, muni de la topologie de la convergence compacte (dite aussi compacte-ouverte), dont les ouverts de base sont les ensembles $\Omega(K, U)$ des $u$ telles que $u(K) \subset U$, pour $K$ compact et $U$ ouvert. La question est de savoir si la bijection envoie $C(E \times F, G)$ sur $C(E, C(F, G))$ et si elle est un homéomorphisme. L'intérêt est concret : une homotopie $E \times [0, 1] \to G$ devient un chemin dans $C(E, G)$, et une famille continue de fonctions dépendant d'un paramètre devient une seule application continue à valeurs dans un espace de fonctions.

\emph{La page.} Le corollaire~2 (PDF, p.~32--33) suppose $F$ localement compact et $E$ séparé, et affirme que la bijection naturelle induit un homéomorphisme de $C_c(E \times F, G)$ sur $C_c(E, C_c(F, G))$. La démonstration a deux temps. La continuité dans le sens direct vient de ce que l'image réciproque de $\Omega(L, \Omega(K, U))$ est $\Omega(L \times K, U)$. Pour la réciproque, on part d'un compact $M$ de $E \times F$ et d'une $f$ qui envoie $M$ dans $U$, et l'on recouvre $M$ par un nombre fini de rectangles compacts contenus dans $f^{-1}(U)$. C'est là que sert l'hypothèse sur $E$ : les projections $K$ et $L$ de $M$ sont compactes, et $K \times L$ est un compact séparé, donc normal, où l'on sait découper. Une note dactylographiée qualifie la preuve de l'ancienne rédaction de \q{canulée} et soupçonne que le corollaire est faux si $E$ n'est pas séparé ; elle en laisse la vérification à Dieudonné, comme exercice.

\emph{Deux conventions.} Chez Bourbaki, un compact est séparé par définition, et un espace localement compact aussi. Dans mathlib, \texttt{IsCompact} ne demande que la propriété de Borel--Lebesgue, et \texttt{LocallyCompactSpace} demande seulement que chaque point ait une base de voisinages compacts, sans séparation. mathlib contient déjà l'homéomorphisme \texttt{Homeomorph.curry}, mais sous l'hypothèse que les deux facteurs $X$ et $Y$ sont localement compacts. Un espace séparé qui n'est pas localement compact, comme $\mathbf Q$, n'y entre pas, alors que la page le couvre.

\begin{definition}
On dit qu'un espace $X$ vérifie (P) si, pour tout compact $K$ de $X$, tout point $x$ de $K$ et tout voisinage $V$ de $x$ dans $X$, il existe un compact $C \subset V$ qui est un voisinage de $x$ dans $K$.
\end{definition}

Tout espace localement compact vérifie (P), en prenant pour $C$ un voisinage compact de $x$ dans $X$. Un espace séparé la vérifie aussi : un compact séparé est localement compact. Plus généralement, il suffit que $X$ soit préregulier (deux points qui n'ont pas les mêmes voisinages ont des voisinages disjoints, \texttt{R1Space}), car un compact préregulier est régulier, donc localement compact.

\begin{enonce}[Corollaire~2, dans les conventions de mathlib]
Soient $X$, $Y$, $Z$ des espaces topologiques. Si $X$ vérifie (P), l'application $f \mapsto (x \mapsto f(x, \cdot))$ de $C(X \times Y, Z)$ dans $C(X, C(Y, Z))$ est continue, sans hypothèse sur $Y$ ni sur $Z$. Si de plus $Y$ est localement compact, c'est un homéomorphisme. En particulier, c'est un homéomorphisme dès que $X$ est séparé et $Y$ localement compact.
\end{enonce}

\begin{proof}
On utilise deux fois le même principe : une propriété des parties d'un compact qui est stable par réunion finie et par passage à une partie, et qui est vraie au voisinage de chaque point, est vraie du compact entier. Pour la continuité directe, on fixe $f$, un compact $K$ de $X$ et un ouvert $W$ de $C(Y, Z)$ contenant les fonctions partielles $f(x, \cdot)$ pour $x \in K$. Il suffit de montrer que chaque point $x$ de $K$ a un voisinage $C$ dans $K$ tel que, pour $f'$ assez proche de $f$, les $f'(x', \cdot)$ avec $x' \in C$ restent dans $W$. Or, si $f(x, \cdot)$ envoie un compact $L$ dans un ouvert $U$, le lemme du tube donne un voisinage $V$ de $x$ tel que $f(V \times L) \subset U$. La propriété (P) fournit un compact $C \subset V$ voisinage de $x$ dans $K$, et l'ouvert $\Omega(C \times L, U)$ contient $f$. Toute $f'$ dans cet ouvert a ses fonctions partielles en les points de $C$ dans $\Omega(L, U)$, et des intersections finies de tels ouverts suffisent pour $W$.

Pour la réciproque, on fixe $g \in C(X, C(Y, Z))$, un compact $M$ de $X \times Y$ et un ouvert $U$ avec $g(x)(y) \in U$ pour $(x, y) \in M$. En un point $(x, y)$ de $M$, la locale compacité de $Y$ donne un voisinage compact $L$ de $y$ avec $g(x)(L) \subset U$, puis la continuité de $g$ un voisinage $V$ de $x$ avec $g(V) \subset \Omega(L, U)$. On applique (P) au compact $p_1(M)$, projection de $M$ sur $X$ : il existe un compact $C \subset V$ voisinage de $x$ dans $p_1(M)$. Le rectangle $C \times L$ est un voisinage de $(x, y)$ dans $M$, et l'ouvert $\Omega(C, \Omega(L, U))$ contient $g$. Le principe de recouvrement fait le reste. Les deux applications sont inverses l'une de l'autre par définition.
\end{proof}

\begin{remarque}
Sous les hypothèses de la page, $E$ séparé et $F$ localement compact séparé, tous les compacts de $E$, de $F$ et de $E \times F$ sont séparés : les deux conventions donnent les mêmes topologies, et l'énoncé démontré contient le corollaire~2. Il le renforce de deux façons : $F$ n'a pas besoin d'être séparé, et $E$ peut être seulement préregulier, ou localement compact au sens de mathlib.
\end{remarque}

\emph{Vérification en Lean.} Lean vérifie l'énoncé ci-dessus dans les conventions de mathlib, avec les deux cas où l'on établit (P) : $X$ localement compact et $X$ préregulier. Il ne dit rien de la question posée par la note.

\emph{La question de la note.} La démonstration n'utilise $E$ qu'à travers (P), au même endroit que la page. Savoir si l'on peut s'en passer reste ouvert ici, dans les deux conventions, et la différence entre elles déplace la difficulté. Dans celle de Bourbaki, la topologie de la convergence compacte ne voit que les compacts séparés, et un compact séparé vérifie toujours (P) en lui-même ; la continuité directe vaut alors pour tout $E$. Pour la réciproque, il faut appliquer (P) à la projection sur $E$ d'un compact séparé $M$ de $E \times F$, et cette projection peut ne pas être séparée quand $E$ ne l'est pas. Sur la droite à deux origines $0'$ et $0''$, par exemple, la partie $\{(0', y_1), (0'', y_2)\}$ avec $y_1 \neq y_2$ est un compact séparé, mais sa projection $\{0', 0''\}$ ne l'est pas. Dans cet exemple, tout se passe bien, car la droite à deux origines est localement compacte au sens de mathlib. Un contre-exemple devrait donc chercher un espace $E$ qui n'est ni préregulier ni localement compact ; aucun n'a été trouvé, et le travail d'Escardó et Heckmann sur les topologies des espaces de fonctions, où la question a sa place, n'a pas été consulté.
